\chapter{TỔNG QUAN}
\label{ch:tong-quan}

Chương này phân tích các hướng nghiên cứu liên quan đến bài toán dự báo xu hướng cổ
phiếu từ dữ liệu giá và tin tức. Việc khảo sát được tổ chức theo ba trụ. Trụ thứ nhất
là phân tích cảm xúc tài chính, gồm các backbone ngôn ngữ và vấn đề thích nghi miền.
Trụ thứ hai là mô hình chuỗi thời gian, từ baseline chỉ giá đến LSTM và Transformer.
Trụ thứ ba là hợp nhất tin tức với giá, trong đó các nghiên cứu được phân biệt theo
điểm cảm xúc, embedding ngữ cảnh, sự kiện và LLM. Sau cùng, chương đối chiếu các hướng
theo dữ liệu, mục tiêu, cách đánh giá và mức độ tái lập để xác định vị trí của đề tài.

Một nguyên tắc quan trọng khi đọc các công trình trước là không đồng nhất kết quả tự
báo cáo với bằng chứng có thể chuyển giao. Phần lớn tài liệu quốc tế sử dụng dữ liệu
tiếng Anh, thị trường Mỹ hoặc Trung Quốc, khung thời gian và cách chia mẫu riêng. Vì
vậy, các công trình được dùng ở đây chủ yếu để rút ra bài học về biểu diễn, kiến trúc
và thiết kế thực nghiệm. Mọi mức cải thiện của hệ thống trên HOSE phải được kiểm tra
lại bằng dữ liệu và protocol của đề tài.

\section{Phân tích cảm xúc tin tức tài chính}
\label{sec:tq-sentiment}

Phân tích cảm xúc tin tức tài chính là bài toán suy ra thái độ hoặc tác động kỳ vọng
của một đoạn văn đối với doanh nghiệp hay thị trường. Trong đề tài này, sentiment
không được hiểu như cảm xúc tâm lý chung của tác giả, mà là tín hiệu định hướng cho
phản ứng kỳ vọng của giá. Một bài viết có thể được gán tích cực nếu nội dung hàm ý
triển vọng tốt hơn cho doanh nghiệp, tiêu cực nếu hàm ý rủi ro hoặc suy giảm, và trung
tính nếu thông tin chủ yếu mô tả sự kiện mà chưa cho thấy hướng tác động rõ ràng. Cách
định nghĩa theo tác động kỳ vọng giúp nhãn phù hợp hơn với mục tiêu dự báo, nhưng cũng
đòi hỏi hướng dẫn gán nhãn chặt chẽ vì nhiều tin có thể chứa cả hai chiều thông tin.

\subsection{Từ mô hình ngôn ngữ tổng quát đến mô hình miền tài chính}

BERT và các biến thể Transformer tạo biểu diễn ngữ cảnh bằng cách cho mỗi token tương
tác với các token khác trong chuỗi. Khi được tinh chỉnh trên dữ liệu có nhãn, backbone
này có thể học những dấu hiệu ngôn ngữ phù hợp với tác vụ phân loại. FinBERT là một
ví dụ về việc đưa mô hình BERT vào miền truyền thông tài chính, qua đó tạo nền tảng
cho các pipeline sentiment rồi dự báo chuyển động cổ phiếu \cite{Yang2020_200608097v2}.
Điểm đáng chú ý là từ điển tài chính không chỉ gồm các từ có sắc thái tích cực hoặc
tiêu cực trong ngôn ngữ thông thường. Một từ như ``tăng'' có thể mô tả chi phí tăng,
doanh thu tăng hoặc rủi ro tăng, nên cần xét ngữ cảnh và chủ thể được đề cập.

Đối với tiếng Việt, PhoBERT là lựa chọn backbone phù hợp vì được tiền huấn luyện trên
ngữ liệu tiếng Việt và có tokenizer tương ứng với đặc trưng tách từ của ngôn ngữ này
\cite{Nguyen2020_200300744v3}. Tuy nhiên, PhoBERT không tự động trở thành mô hình
sentiment tài chính chỉ nhờ tiền huấn luyện. Mô hình vẫn cần được tinh chỉnh trên dữ
liệu có nhãn của miền tài chính, cần kiểm tra sự khác biệt giữa tiêu đề và nội dung,
và cần đánh giá trên phần dữ liệu theo thời gian không trùng với dữ liệu huấn luyện.

\subsection{Các công trình sentiment và chuyển động cổ phiếu quốc tế}

Một dòng nghiên cứu kết hợp trực tiếp sentiment với mô hình chuỗi thời gian là
FinBERT-LSTM. Halder xây dựng pipeline trong đó tín hiệu từ tin tức được đưa vào LSTM
để dự báo giá, còn Gu và cộng sự phát triển một biến thể cùng dòng với cách tổ chức
dữ liệu và mô hình được cải tiến hơn \cite{Halder2022_221107392v1,
Gu2024_240716150v1}. Các công trình này có giá trị như mẫu kiến trúc cho việc tách
nhánh văn bản và nhánh giá, nhưng có những khác biệt cần lưu ý: dữ liệu chủ yếu là
tiếng Anh, mục tiêu là giá hoặc biến động trong bối cảnh khác, và cách ghép đặc trưng
không thể được xem là giao thức chuẩn cho HOSE.

Chen đề xuất sử dụng embedding ngữ cảnh từ BERT thay vì chỉ nén mỗi bài viết thành một
điểm sentiment, sau đó đưa biểu diễn vào mô hình RNN cho dự báo chuyển động. Công trình
này cho thấy việc giữ thông tin ngữ cảnh có thể là một hướng mở rộng đáng chú ý, đồng
thời nhấn mạnh việc đánh giá bằng mô phỏng giao dịch thay vì chỉ nhìn vào accuracy
\cite{Chen2021_210708721v1}. Tuy nhiên, mô phỏng giao dịch vẫn phụ thuộc vào giả định
về chi phí, thời điểm khớp lệnh và khả năng thực hiện, nên không thể thay thế việc
kiểm tra phân loại ngoài mẫu trong graduation floor của đề tài.

Jiang nghiên cứu sentiment tài chính bằng FinBERT gắn với dự báo chuyển động, còn Siala
và cộng sự tập trung trực tiếp vào câu hỏi sentiment do LLM tạo ra cải thiện dự báo ở
mức nào \cite{Jiang2023_230602136v3,Siala2026_260200086v3}. Những hướng này củng cố
quan điểm rằng biến sentiment nên được đánh giá bằng phép so sánh có và không có tín
hiệu trên cùng dữ liệu. Chúng không cho phép kết luận trước rằng tín hiệu do PhoBERT
tạo ra sẽ có tác dụng tương tự trên văn bản tiếng Việt.

\subsection{Đặc thù của sentiment tài chính tiếng Việt}

Trong bối cảnh tiếng Việt, khó khăn đầu tiên là thiếu dữ liệu có nhãn lớn, nhất quán
và có mô tả rõ ràng về tiêu chí gán nhãn. Một bộ dữ liệu tiêu đề có thể thuận tiện để
khởi tạo fine-tune, nhưng tiêu đề ngắn thường bỏ qua điều kiện, thời hạn và chủ thể
được mô tả trong thân bài. Ngược lại, dùng toàn văn giúp tăng ngữ cảnh nhưng làm tăng
chi phí làm sạch, giới hạn độ dài và rủi ro một bài chứa nhiều sự kiện trái chiều.

Khó khăn thứ hai là sự trùng lặp theo nguồn. Cùng một thông cáo có thể được đăng lại
ở nhiều trang hoặc gắn với nhiều mã. Nếu một bản sao xuất hiện ở cả train và test,
mô hình có thể ghi nhớ câu chữ thay vì học quy luật sentiment. Vì thế, loại trùng cần
được thực hiện ở cấp văn bản chuẩn hóa, đồng thời ghi lại URL, nguồn, thời điểm thu
thập và hash dữ liệu.

Khó khăn thứ ba là sentiment của một bài không nhất thiết có tác động cùng chiều với
lợi suất. Tin có thể đã được thị trường biết trước, có thể mô tả một sự kiện dài hạn
nhưng không ảnh hưởng phiên kế tiếp, hoặc chỉ tác động đến biến động chứ không tác
động đến chiều giá. Tài liệu nguồn của đề tài cũng xem việc phản ứng giá yếu hoặc
không ổn định là khả năng cần kiểm tra, không phải dấu hiệu thất bại của mô hình.

\subsection{Đánh giá chất lượng nhãn và đầu ra sentiment}

Đánh giá sentiment cần tách khỏi đánh giá dự báo. Ở tầng thứ nhất, mô hình được kiểm
tra bằng accuracy, precision, recall và F1 theo lớp, trong đó macro-F1 được ưu tiên khi
phân bố lớp không cân bằng. Tập in-domain cần được gán nhãn theo hướng dẫn cố định;
độ đồng thuận giữa hai người gán có thể báo cáo bằng Cohen's kappa, còn từ ba người
gán trở lên có thể dùng Fleiss' kappa. Các mẫu bất đồng phải được ghi nhận thay vì
loại bỏ tùy ý cho đến khi đạt điểm số mong muốn.

Đầu ra được dùng ở tầng dự báo không chỉ là nhãn argmax. Với mỗi bài viết, mô hình
trả về vector xác suất ba chiều
\[
  \mathbf{p}_{i,j,t} = (p_{\mathrm{NEG}},p_{\mathrm{NEU}},p_{\mathrm{POS}}),
  \qquad \sum_{c} p_{c}=1,
\]
trong đó $i$ là mã cổ phiếu, $j$ là bài viết và $t$ là ngày được ánh xạ. Vector này
giữ lại mức độ không chắc chắn của mô hình. Sau đó, các vector được tổng hợp theo mã
và ngày cùng với số tin, độ phân tán và cờ có tin. Cách làm này cho phép mô hình dự
báo phân biệt ngày có nhiều tin trung tính với ngày hoàn toàn không có tin.

\section{Dự báo xu hướng giá bằng mô hình chuỗi thời gian}
\label{sec:tq-timeseries}

Nhánh chuỗi thời gian biểu diễn trạng thái thị trường bằng lịch sử OHLCV, lợi suất và
các chỉ báo kỹ thuật. Mục tiêu của nhánh này không phải tái tạo chính xác mức giá,
mà là cung cấp một baseline có năng lực để kiểm tra sentiment có bổ sung thông tin
hay không. Các nghiên cứu dự báo xu hướng quốc tế thường sử dụng lịch sử giá kết hợp
đặc trưng kỹ thuật, sau đó mở rộng sang mô hình sâu hoặc mô hình quan hệ giữa nhiều
mã \cite{Feng2018_181009936v2,Xu2021_211013716v2}.

\subsection{Baseline chỉ giá và các mô hình cổ điển}

Baseline majority dự đoán lớp xuất hiện nhiều nhất trong phần huấn luyện. Baseline
random sinh dự đoán theo phân bố lớp của phần huấn luyện và dùng seed cố định. Hai
baseline này có vai trò kiểm tra xem mô hình học được tín hiệu vượt qua một quy tắc
đơn giản hay không. Nếu một mô hình sâu chỉ nhỉnh hơn majority ở accuracy nhưng thua
macro-F1, kết luận về hiệu quả cần dựa vào chỉ số cân bằng hơn.

Nhóm mô hình cổ điển chỉ giá gồm Logistic Regression, Random Forest và XGBoost. Chúng
được huấn luyện trên cùng tập đặc trưng kỹ thuật với dữ liệu được chuẩn hóa trong
từng cửa sổ. Mục đích không phải chạy thật nhiều thuật toán, mà tạo một mốc kiểm tra
đủ rõ cho nhánh giá. Nếu mô hình có sentiment không vượt được mốc này, việc thêm một
kiến trúc sâu hơn chưa chắc giải quyết đúng vấn đề dữ liệu.

\subsection{LSTM cho phụ thuộc ngắn hạn và dài hạn}

LSTM sử dụng trạng thái ẩn và trạng thái ô nhớ để cập nhật thông tin qua các thời
điểm. Nhờ các cổng vào, cổng quên và cổng ra, LSTM có thể giữ hoặc loại bỏ một phần
thông tin trong cửa sổ lịch sử. Trong các pipeline dự báo tài chính gắn với tin tức,
LSTM thường được dùng như bộ mã hóa chuỗi hoặc bộ dự báo cuối
\cite{Halder2022_221107392v1,Gu2024_240716150v1}. Các mốc accuracy cao được tài liệu
nguồn tổng hợp cho dữ liệu Việt Nam phải được đọc thận trọng nếu không có kiểm tra
leakage, metric theo lớp hoặc metric giao dịch. Vì khác biệt về thị trường, cách tạo
nhãn và cách chia dữ liệu, các mốc đó chỉ là bối cảnh tham khảo, không phải kết quả của
đề tài.

Ở bình diện quốc tế, các nghiên cứu dự báo chuyển động cũng dùng nhiều biến thể mạng
hồi quy và thường kết hợp thêm thông tin từ tin tức hoặc quan hệ giữa các mã
\cite{Halder2022_221107392v1,Chen2021_210708721v1}. Vì mục tiêu của đề tài là đo
incremental value của sentiment, LSTM chỉ giá và LSTM hai nhánh được giữ cùng cửa sổ,
chiến lược chia dữ liệu và ngân sách huấn luyện. Đây là điều kiện cần để chênh lệch
giữa hai mô hình có thể được quy về việc thêm modality, thay vì quy về thay đổi quy mô
mạng.

\subsection{Transformer cho chuỗi thời gian}

Transformer dùng self-attention để tính mức liên hệ giữa các vị trí trong chuỗi. Khác
với việc chỉ truyền trạng thái tuần tự, attention có thể kết nối trực tiếp các thời
điểm xa nhau và xử lý nhiều kênh đặc trưng trong cùng một lớp. Trong dự báo cổ phiếu,
Transformer còn có thể được mở rộng để mô hình hóa phụ thuộc giữa nhiều mã hoặc giữa
nhiều loại tín hiệu. Boyle và Kalita nghiên cứu hướng spatiotemporal Transformer cho
dự báo chuyển động, còn Omranpour và cộng sự nghiên cứu Transformer bậc cao cho dữ
liệu chuỗi thời gian đa phương thức \cite{Boyle2023_230503835v1,
Omranpour2024_241210540v1}.

Temporal Fusion Transformer được xem trong đề tài như tầng mở rộng của thang mô hình.
Ý tưởng chính là kết hợp xử lý chuỗi cục bộ, cơ chế chọn biến, các cổng điều tiết và
attention để mô hình có thể sử dụng nhiều nhóm biến. Với bài toán hiện tại, TFT phải
được cấu hình cho phân loại ba lớp, sử dụng đầu ra kích thước 3 và loss phân loại.
Không được lấy cấu hình hồi quy phân vị mặc định rồi diễn giải kết quả như phân loại.
Các trọng số attention hoặc variable importance của TFT chỉ là bằng chứng mô tả phụ
thuộc vào mô hình, không phải bằng chứng nhân quả.

\subsection{Dịch chuyển phân phối và quan hệ giữa các mã}

Một số hướng nghiên cứu quốc tế không chỉ mô hình hóa từng mã độc lập. HIST khai thác
thông tin chia sẻ giữa mã và khái niệm, REST đưa sự kiện vào quan hệ giữa các cổ phiếu,
còn DoubleAdapt tập trung vào thích nghi khi phân phối dữ liệu thay đổi
\cite{Xu2021_211013716v2,Xu2021_210207372v2,Zhao2023_230609862v3}. MTMD, DishFT-GNN
và S$^{3}$G tiếp tục phát triển các dạng bộ nhớ đa quy mô, đồ thị động hoặc biểu diễn
trạng thái để xử lý nhiễu và quan hệ liên mã \cite{Wang2022_221208656v2,
Liu2025_250210776v1,Lu2026_260324236v1}.

Các công trình này cho thấy bài toán dự báo xu hướng có thể chịu ảnh hưởng của quan
hệ ngành, khái niệm và concept drift. Tuy nhiên, đưa đồ thị nhiều mã vào graduation
floor sẽ làm tăng số giả định về quan hệ liên cổ phiếu và độ phức tạp của đánh giá.
Đề tài trước hết giữ mô hình hai nhánh trên từng mã, dùng nhiều mã trong cùng lịch
walk-forward để kiểm tra tính ổn định. Quan hệ liên mã và thích nghi online được dành
cho các hướng phát triển sau khi pipeline cơ sở đã được khóa.

\section{Kết hợp tin tức và giá (multimodal)}
\label{sec:tq-fusion}

Hợp nhất multimodal là quá trình đưa thông tin từ văn bản và chuỗi số vào một bộ dự
báo chung. Hai nguồn khác nhau ở kích thước, tần suất, độ trễ và mức độ tin cậy, nên
phép nối đặc trưng đơn giản có thể là một baseline hữu ích nhưng chưa chắc mô hình hóa
được tương tác giữa chúng. Các nghiên cứu quốc tế cho thấy có thể phân biệt một số
mức hợp nhất, từ điểm sentiment theo ngày, embedding văn bản, sự kiện chi tiết đến
LLM và cross-attention \cite{Halder2022_221107392v1,Elahi2024_241101368v1}.

\subsection{Hợp nhất bằng điểm sentiment}

Ở mức đơn giản nhất, mô hình NLP biến mỗi bài thành nhãn hoặc vector xác suất. Các
vector được gom theo mã và ngày, sau đó nối với vector kỹ thuật trước khi đưa vào
Logistic Regression, cây quyết định hoặc LSTM. Ưu điểm của cách này là dễ kiểm tra,
dễ ablation và ít tốn tài nguyên. Nhược điểm là toàn bộ thông tin trong nhiều bài
được nén vào một số trung bình, làm mất thứ tự bài, sự kiện và quan hệ giữa câu chữ
với đặc trưng giá.

FinBERT-LSTM là mẫu điển hình của cách tiếp cận sentiment rồi chuỗi thời gian
\cite{Halder2022_221107392v1,Gu2024_240716150v1}. Với đề tài, cách tiếp cận này được
chuyển sang PhoBERT và dữ liệu tiếng Việt. Tuy nhiên, bản thân việc dùng PhoBERT không
đủ để tạo đóng góp. Đóng góp cần được chứng minh bằng việc mô hình có sentiment được
đặt cạnh mô hình price-only dưới cùng protocol, và bằng việc kiểm tra các trường hợp
không có tin.

\subsection{Hợp nhất bằng embedding ngữ cảnh và sự kiện}

Thay vì chỉ dùng điểm sentiment, một hướng khác giữ lại embedding ngữ cảnh của tiêu
đề hoặc bài viết. Chen sử dụng contextualized embedding từ BERT cho dự báo movement,
cho thấy cách biểu diễn này có thể chứa các tín hiệu vượt ra ngoài phân cực ba lớp
\cite{Chen2021_210708721v1}. Một hướng giàu hơn nữa là trích xuất sự kiện chi tiết,
chẳng hạn chủ thể, loại sự kiện, hướng tác động và thời gian. Chen và cộng sự nghiên
cứu việc đưa các sự kiện fine-grained vào dự báo chuyển động, qua đó gợi ý rằng cảm
xúc chỉ là một phần của thông tin văn bản \cite{Chen2019_191005078v1}.

Tuy nhiên, embedding toàn văn có kích thước lớn và có thể chứa thông tin không liên
quan. Event extraction cũng cần schema, gán nhãn và kiểm tra mapping phức tạp. Vì đề
tài đặt mục tiêu tốt nghiệp có thể tái lập, graduation floor chọn xác suất sentiment
được tổng hợp theo ngày. Embedding và event-level signal được xem là phép thử nhạy
cảm hoặc hướng phát triển, không thay thế baseline chính.

\subsection{Hợp nhất động và LLM}

Các công trình gần đây dùng LLM để kết hợp dữ liệu giá có cấu trúc với tin tức hoặc
để mô hình hóa cách tin được lan truyền. Elahi và Taghvaei đặt bài toán kết hợp dữ
liệu tài chính với bài viết qua LLM, còn Liang và cộng sự đưa yếu tố dissemination và
ngữ cảnh vào FinGPT \cite{Elahi2024_241101368v1,Liang2024_241210823v2}. Những hướng
này có thể giữ nhiều ngữ cảnh hơn, nhưng phụ thuộc vào kích thước mô hình, chi phí,
phiên bản checkpoint và khả năng tái lập. Nếu dùng mô hình thương mại, việc xác định
đúng thời điểm kiến thức và tái tạo cùng đầu ra còn khó hơn.

Với bài toán của đề tài, PhoBERT mở được chọn thay cho LLM thương mại vì có thể ghi
nhận checkpoint, tokenizer, seed và cấu hình. Lựa chọn này không có nghĩa PhoBERT luôn
tốt hơn LLM, mà giúp phạm vi nghiên cứu có thể kiểm tra trong điều kiện tài nguyên của
đồ án. Siala và cộng sự cho thấy câu hỏi ``sentiment có thật sự làm tăng chất lượng
dự báo'' cần được đo riêng, một điểm phù hợp với thiết kế ablation của đề tài
\cite{Siala2026_260200086v3}.

\subsection{Các cơ chế fusion}

Có thể phân loại cơ chế fusion thành ba nhóm:

\begin{enumerate}[nosep]
  \item \textbf{Early fusion}: chuẩn hóa hai nhóm đặc trưng rồi nối chúng trước bộ
    mã hóa. Cách này đơn giản nhưng dễ để nhóm có nhiều chiều lấn át nhóm còn lại.
  \item \textbf{Late fusion}: mỗi nhánh tạo ra biểu diễn hoặc dự đoán riêng, sau đó
    một lớp kết hợp các đầu ra. Cách này tách biệt hơn nhưng có thể bỏ lỡ tương tác
    chi tiết giữa một sự kiện và một trạng thái giá.
  \item \textbf{Gated hoặc cross-attention fusion}: một nhánh học cách chọn thông tin
    từ nhánh kia theo thời điểm và ngữ cảnh. Cách này giàu biểu đạt hơn, nhưng cần
    nhiều dữ liệu và dễ làm tăng phương sai khi số mã chỉ ở mức 10.
\end{enumerate}

Đề tài dùng concat sau hai nhánh LSTM làm mức sàn vì đây là phép hợp nhất có thể kiểm
tra, dễ giữ ngân sách tham số gần với price-only LSTM và phù hợp với ablation. Gated
cross-attention hoặc dual-stream Transformer được nêu ở Chương~\ref{ch:huong-phat-trien},
không được mặc định xem là đối thủ bắt buộc của graduation floor. Higher Order
Transformers là ví dụ cho việc đưa nhiều chuỗi vào một Transformer, nhưng mức phức tạp
của hướng đó phù hợp hơn với tầng mở rộng sau khi baseline đã ổn định
\cite{Omranpour2024_241210540v1}.

\subsection{Đánh giá multimodal và nguy cơ diễn giải quá mức}

Một mô hình multimodal có thể đạt điểm cao vì nhiều lý do khác nhau. Tín hiệu sentiment
có thể phản ánh lại các biến giá đã được dùng trong cùng ngày, tin có thể xuất hiện
sau thời điểm quyết định, hoặc dữ liệu train và test có thể chứa các bản tin trùng
lặp. Vì vậy, chỉ so sánh một mô hình fusion với một mô hình price-only khác cấu hình
không đủ để khẳng định giá trị của văn bản.

Thiết kế của đề tài yêu cầu bốn điều kiện. Thứ nhất, mapping tin theo timestamp phải
được audit. Thứ hai, cùng một cửa sổ, seed, preprocessing và holdout phải được dùng
cho hai nhánh ablation. Thứ ba, chỉ số chính phải tính trên dự đoán ngoài mẫu theo
thứ tự thời gian. Thứ tư, chênh lệch phải đi kèm độ bất định, chẳng hạn bootstrap theo
khối ngày. Cách đánh giá giàu hơn như mô phỏng giao dịch có thể bổ sung trong tương
lai, như tinh thần được minh họa ở các nghiên cứu dùng trading simulation
\cite{Chen2021_210708721v1,Feng2018_181009936v2}.

\section{Khoảng trống nghiên cứu và định vị đề tài}
\label{sec:tq-gap}

\subsection{Tổng hợp theo ba trụ}

Từ khảo sát, khoảng trống không được hiểu là chưa từng có ai dùng sentiment, LSTM hay
Transformer riêng lẻ. Các thành phần cơ bản đã xuất hiện trong nhiều dòng nghiên cứu:
FinBERT cho sentiment tài chính tiếng Anh \cite{Yang2020_200608097v2}, PhoBERT cho
tiếng Việt \cite{Nguyen2020_200300744v3}, FinBERT-LSTM cho pipeline tin tức và giá
\cite{Halder2022_221107392v1,Gu2024_240716150v1}, contextualized embedding cho
movement \cite{Chen2021_210708721v1}, và Transformer cho chuỗi thời gian đa phương
thức hoặc spatiotemporal \cite{Omranpour2024_241210540v1,Boyle2023_230503835v1}.

Khoảng trống nằm ở sự kết hợp và cách kiểm chứng trong bối cảnh cụ thể. Trong phạm vi
tài liệu được khảo sát cho đề tài, chưa có một pipeline được khóa đầy đủ các thành
phần sau trên cùng rổ nhiều mã HOSE: backbone PhoBERT mở cho tin tiếng Việt, ánh xạ
point-in-time theo cutoff, nhánh giá với đặc trưng kỹ thuật, mô hình hai nhánh, ablation
có và không có sentiment, walk-forward không rò rỉ và báo cáo chênh lệch có khoảng tin
cậy. Đây là tuyên bố về phạm vi khảo sát và thiết kế của đề tài, không phải tuyên bố
rằng không tồn tại bất kỳ triển khai nào khác.

\subsection{So sánh các hướng nghiên cứu}

\begin{table}[H]
  \centering
  \caption{So sánh các nhóm nghiên cứu liên quan với phạm vi đề tài.}
  \label{tab:tq-so-sanh}
  \begin{tabular}{p{3.0cm} p{4.1cm} p{4.1cm} p{2.6cm}}
    \toprule
    \textbf{Nhóm} & \textbf{Đặc điểm chính} & \textbf{Hạn chế khi áp dụng cho đề tài} & \textbf{Vai trò} \\
    \midrule
    Sentiment tài chính & FinBERT và PhoBERT tạo nhãn hoặc xác suất cảm xúc cho văn bản
    \cite{Yang2020_200608097v2,Nguyen2020_200300744v3} & Khác ngôn ngữ, miền,
    cách gán nhãn và thời điểm tin & Backbone, tầng tiền xử lý \\
    Price-only LSTM & Mô hình hóa cửa sổ OHLCV và chỉ báo kỹ thuật theo thời gian
    \cite{Halder2022_221107392v1} & Không kiểm tra giá trị gia tăng của văn bản & Baseline
    sâu \\
    Transformer movement & Attention cho phụ thuộc thời gian, đa phương thức hoặc
    liên mã \cite{Boyle2023_230503835v1,Omranpour2024_241210540v1} & Phức tạp,
    chưa có bằng chứng trực tiếp trên HOSE & TFT stretch \\
    News-price fusion & Nối sentiment, embedding, event hoặc LLM với giá
    \cite{Chen2021_210708721v1,Elahi2024_241101368v1} & Dễ trộn lẫn khác biệt dữ liệu,
    rò rỉ và giả định thị trường & Mẫu kiến trúc \\
    Graph và meta-learning & Khai thác quan hệ liên mã và concept drift
    \cite{Xu2021_211013716v2,Zhao2023_230609862v3} & Ngoài phạm vi 10 mã độc lập và
    ngân sách đồ án & Bối cảnh tương lai \\
    \bottomrule
  \end{tabular}
\end{table}

\subsection{Định vị và giả thuyết kiểm chứng}

Đề tài được định vị như một nghiên cứu thực nghiệm multimodal tiếng Việt, không phải
một cuộc thi chọn kiến trúc mới. Mức sàn gồm majority, random, price-only classical,
price-only LSTM và LSTM hai nhánh với concat. TFT chỉ được chạy sau khi mức sàn có số
liệu và đã kiểm tra loss phân loại, output kích thước 3 cùng holdout. Cách xếp tầng
này kế thừa tinh thần xây dựng đối chứng và kiểm tra đặc trưng trước khi tăng độ phức
tạp, phù hợp với việc các nghiên cứu price-only cho thấy tiền xử lý và lựa chọn đặc
trưng có thể ảnh hưởng lớn đến kết quả \cite{Feng2018_181009936v2}.

Giả thuyết chính là vector xác suất sentiment tổng hợp có thể cung cấp thông tin gia
tăng cho dự báo xu hướng phiên kế tiếp sau khi đã kiểm soát lịch sử giá. Giả thuyết
phụ là phần gia tăng, nếu tồn tại, không nhất thiết ổn định ở mọi ngày mà có thể rõ
hơn quanh ngày có nhiều tin hoặc giai đoạn biến động cao. Đây là giả thuyết mô tả cần
được kiểm tra bằng phân tầng dữ liệu; không được biến nó thành khẳng định rằng tin
luôn gây ra biến động.

\subsection{Tiêu chí để tuyên bố đóng góp}

Một tuyên bố đóng góp chỉ được đưa ra khi thỏa mãn các điều kiện:

\begin{itemize}[nosep]
  \item dữ liệu và phạm vi mã được mô tả rõ, bao gồm bản ghi thiếu hoặc bị loại;
  \item cutoff, công thức nhãn, phép chia theo thời gian và quy trình fit tham số được
    ghi trong manifest;
  \item kết quả sentiment và kết quả dự báo được tách thành hai tầng;
  \item mô hình có sentiment và price-only dùng cùng protocol, seed hoặc tập seed đã
    định trước;
  \item macro-F1, balanced accuracy, macro OvR-AUC và chỉ số theo lớp được báo cáo
    cùng khoảng tin cậy khi có thể;
  \item các kết luận nêu rõ tính tự báo cáo của tài liệu tham khảo và không suy rộng
    ngoài 10 mã thuộc VN30 trong giai đoạn nghiên cứu.
\end{itemize}

Theo đó, đóng góp an toàn nhất của đề tài là một pipeline tái lập để đo lường phần
đóng góp biên của sentiment tiếng Việt dưới đánh giá walk-forward có kiểm soát. Nếu
ablation không cho thấy cải thiện, kết quả đó vẫn trả lời được câu hỏi nghiên cứu và
giúp tránh kết luận quá mức từ một mô hình phức tạp.
